% Options for packages loaded elsewhere
\PassOptionsToPackage{unicode}{hyperref}
\PassOptionsToPackage{hyphens}{url}
\documentclass[
  ignorenonframetext,
]{beamer}
\newif\ifbibliography
\usepackage{pgfpages}
\setbeamertemplate{caption}[numbered]
\setbeamertemplate{caption label separator}{: }
\setbeamercolor{caption name}{fg=normal text.fg}
\beamertemplatenavigationsymbolsempty
% remove section numbering
\setbeamertemplate{part page}{
  \centering
  \begin{beamercolorbox}[sep=16pt,center]{part title}
    \usebeamerfont{part title}\insertpart\par
  \end{beamercolorbox}
}
\setbeamertemplate{section page}{
  \centering
  \begin{beamercolorbox}[sep=12pt,center]{section title}
    \usebeamerfont{section title}\insertsection\par
  \end{beamercolorbox}
}
\setbeamertemplate{subsection page}{
  \centering
  \begin{beamercolorbox}[sep=8pt,center]{subsection title}
    \usebeamerfont{subsection title}\insertsubsection\par
  \end{beamercolorbox}
}
% Prevent slide breaks in the middle of a paragraph
\widowpenalties 1 10000
\raggedbottom
\AtBeginPart{
  \frame{\partpage}
}
\AtBeginSection{
  \ifbibliography
  \else
    \frame{\sectionpage}
  \fi
}
\AtBeginSubsection{
  \frame{\subsectionpage}
}
\usepackage{iftex}
\ifPDFTeX
  \usepackage[T1]{fontenc}
  \usepackage[utf8]{inputenc}
  \usepackage{textcomp} % provide euro and other symbols
\else % if luatex or xetex
  \usepackage{unicode-math} % this also loads fontspec
  \defaultfontfeatures{Scale=MatchLowercase}
  \defaultfontfeatures[\rmfamily]{Ligatures=TeX,Scale=1}
\fi
\usepackage{lmodern}
\ifPDFTeX\else
  % xetex/luatex font selection
\fi
% Use upquote if available, for straight quotes in verbatim environments
\IfFileExists{upquote.sty}{\usepackage{upquote}}{}
\IfFileExists{microtype.sty}{% use microtype if available
  \usepackage[]{microtype}
  \UseMicrotypeSet[protrusion]{basicmath} % disable protrusion for tt fonts
}{}
\makeatletter
\@ifundefined{KOMAClassName}{% if non-KOMA class
  \IfFileExists{parskip.sty}{%
    \usepackage{parskip}
  }{% else
    \setlength{\parindent}{0pt}
    \setlength{\parskip}{6pt plus 2pt minus 1pt}}
}{% if KOMA class
  \KOMAoptions{parskip=half}}
\makeatother
\usepackage{color}
\usepackage{fancyvrb}
\newcommand{\VerbBar}{|}
\newcommand{\VERB}{\Verb[commandchars=\\\{\}]}
\DefineVerbatimEnvironment{Highlighting}{Verbatim}{commandchars=\\\{\}}
% Add ',fontsize=\small' for more characters per line
\newenvironment{Shaded}{}{}
\newcommand{\AlertTok}[1]{\textcolor[rgb]{1.00,0.00,0.00}{\textbf{#1}}}
\newcommand{\AnnotationTok}[1]{\textcolor[rgb]{0.38,0.63,0.69}{\textbf{\textit{#1}}}}
\newcommand{\AttributeTok}[1]{\textcolor[rgb]{0.49,0.56,0.16}{#1}}
\newcommand{\BaseNTok}[1]{\textcolor[rgb]{0.25,0.63,0.44}{#1}}
\newcommand{\BuiltInTok}[1]{\textcolor[rgb]{0.00,0.50,0.00}{#1}}
\newcommand{\CharTok}[1]{\textcolor[rgb]{0.25,0.44,0.63}{#1}}
\newcommand{\CommentTok}[1]{\textcolor[rgb]{0.38,0.63,0.69}{\textit{#1}}}
\newcommand{\CommentVarTok}[1]{\textcolor[rgb]{0.38,0.63,0.69}{\textbf{\textit{#1}}}}
\newcommand{\ConstantTok}[1]{\textcolor[rgb]{0.53,0.00,0.00}{#1}}
\newcommand{\ControlFlowTok}[1]{\textcolor[rgb]{0.00,0.44,0.13}{\textbf{#1}}}
\newcommand{\DataTypeTok}[1]{\textcolor[rgb]{0.56,0.13,0.00}{#1}}
\newcommand{\DecValTok}[1]{\textcolor[rgb]{0.25,0.63,0.44}{#1}}
\newcommand{\DocumentationTok}[1]{\textcolor[rgb]{0.73,0.13,0.13}{\textit{#1}}}
\newcommand{\ErrorTok}[1]{\textcolor[rgb]{1.00,0.00,0.00}{\textbf{#1}}}
\newcommand{\ExtensionTok}[1]{#1}
\newcommand{\FloatTok}[1]{\textcolor[rgb]{0.25,0.63,0.44}{#1}}
\newcommand{\FunctionTok}[1]{\textcolor[rgb]{0.02,0.16,0.49}{#1}}
\newcommand{\ImportTok}[1]{\textcolor[rgb]{0.00,0.50,0.00}{\textbf{#1}}}
\newcommand{\InformationTok}[1]{\textcolor[rgb]{0.38,0.63,0.69}{\textbf{\textit{#1}}}}
\newcommand{\KeywordTok}[1]{\textcolor[rgb]{0.00,0.44,0.13}{\textbf{#1}}}
\newcommand{\NormalTok}[1]{#1}
\newcommand{\OperatorTok}[1]{\textcolor[rgb]{0.40,0.40,0.40}{#1}}
\newcommand{\OtherTok}[1]{\textcolor[rgb]{0.00,0.44,0.13}{#1}}
\newcommand{\PreprocessorTok}[1]{\textcolor[rgb]{0.74,0.48,0.00}{#1}}
\newcommand{\RegionMarkerTok}[1]{#1}
\newcommand{\SpecialCharTok}[1]{\textcolor[rgb]{0.25,0.44,0.63}{#1}}
\newcommand{\SpecialStringTok}[1]{\textcolor[rgb]{0.73,0.40,0.53}{#1}}
\newcommand{\StringTok}[1]{\textcolor[rgb]{0.25,0.44,0.63}{#1}}
\newcommand{\VariableTok}[1]{\textcolor[rgb]{0.10,0.09,0.49}{#1}}
\newcommand{\VerbatimStringTok}[1]{\textcolor[rgb]{0.25,0.44,0.63}{#1}}
\newcommand{\WarningTok}[1]{\textcolor[rgb]{0.38,0.63,0.69}{\textbf{\textit{#1}}}}
\usepackage{longtable,booktabs,array}
\newcounter{none} % for unnumbered tables
\usepackage{calc} % for calculating minipage widths
\usepackage{caption}
% Make caption package work with longtable
\makeatletter
\def\fnum@table{\tablename~\thetable}
\makeatother
\setlength{\emergencystretch}{3em} % prevent overfull lines
\providecommand{\tightlist}{%
  \setlength{\itemsep}{0pt}\setlength{\parskip}{0pt}}
% Auto-generated by infrastructure/rendering/_pdf_latex_helpers.py::extract_math_font_preamble.
% Minimal header passed to Pandoc via -H so Beamer slides pick up the same math font
% as the combined-PDF path. Adjust the manuscript's preamble.md to change the font globally.
\usepackage{unicode-math}
\setmathfont{latinmodern-math.otf}

% Slide layout defaults for warning-clean scientific decks.
\usepackage{etoolbox}
\IfFileExists{xurl.sty}{\usepackage{xurl}}{}
\IfFileExists{seqsplit.sty}{\usepackage{seqsplit}}{\newcommand{\seqsplit}[1]{#1}}
\protected\def\breakseq#1{\seqsplit{#1}}
\protected\def\breaktt#1{\begingroup\ttfamily\seqsplit{#1}\endgroup}
\setlength{\emergencystretch}{6em}
\tolerance=5000
\hbadness=10000
\hfuzz=1pt
\setlength{\tabcolsep}{2pt}
\AtBeginEnvironment{longtable}{\tiny\renewcommand{\arraystretch}{0.86}\setlength{\tabcolsep}{1pt}}
\AtBeginEnvironment{tabular}{\tiny\renewcommand{\arraystretch}{0.86}\setlength{\tabcolsep}{1pt}}

% Natbib and cross-reference fallbacks — slides don't load natbib
% or cleveref, but combined-PDF manuscript prose may emit these
% commands. The fallback renders citations as a bracketed key list
% and cross-references as detokenized labels so slides don't fail on
% undefined control sequences or raw underscores. \providecommand is
% a no-op if packages load later.
\providecommand{\citep}[1]{[#1]}
\providecommand{\citet}[1]{#1}
\providecommand{\citealp}[1]{#1}
\providecommand{\citeauthor}[1]{#1}
\providecommand{\citeyear}[1]{#1}
\providecommand{\cref}[1]{\texttt{\detokenize{#1}}}
\providecommand{\Cref}[1]{\texttt{\detokenize{#1}}}
\usepackage{bookmark}
\IfFileExists{xurl.sty}{\usepackage{xurl}}{} % add URL line breaks if available
\urlstyle{same}
% fallback for those not using the hyperref driver hyperxmp:
\makeatletter
\@ifundefined{xmpquote}{\newcommand{\xmpquote}[1]{#1}}{}
\makeatother
\hypersetup{
  hidelinks,
  pdfcreator={LaTeX via pandoc}}

\author{\texorpdfstring{}{}}
\date{}

\begin{document}

\section{Reproducibility}\label{sec:reproducibility}

\begin{frame}[allowframebreaks]{Reproducibility}
This section explains how to regenerate every artifact in the study from
a clean checkout. The exemplar's reproducibility guarantee is
structural: each result is produced by a tested function and a thin
script, then injected into the manuscript by generated-variable
substitution --- never transcribed by hand.
\end{frame}

\begin{frame}[fragile,allowframebreaks]{How to regenerate everything}
\protect\phantomsection\label{how-to-regenerate-everything}
From the repository root:

\begin{Shaded}
\begin{Highlighting}[]
\CommentTok{\# 1. Run the analysis (compiles methods, exports artifacts, writes figure + reports)}
\ExtensionTok{uv}\NormalTok{ run python projects/templates/template\_methods\_paper/scripts/methods\_analysis.py}

\CommentTok{\# 2. Run the test suite with the coverage gate}
\ExtensionTok{uv}\NormalTok{ run pytest projects/templates/template\_methods\_paper/tests }\DataTypeTok{\textbackslash{}}
    \AttributeTok{{-}{-}cov}\OperatorTok{=}\NormalTok{projects/templates/template\_methods\_paper/src }\AttributeTok{{-}{-}cov{-}fail{-}under}\OperatorTok{=}\NormalTok{90}

\CommentTok{\# 3. Generate and inject manuscript variables}
\ExtensionTok{uv}\NormalTok{ run python projects/templates/template\_methods\_paper/scripts/z\_generate\_manuscript\_variables.py}

\CommentTok{\# 4. Render the manuscript}
\ExtensionTok{uv}\NormalTok{ run python scripts/03\_render\_pdf.py }\AttributeTok{{-}{-}project}\NormalTok{ templates/template\_methods\_paper}
\end{Highlighting}
\end{Shaded}

Or, end to end via the orchestrated pipeline:

\begin{Shaded}
\begin{Highlighting}[]
\ExtensionTok{uv}\NormalTok{ run python scripts/execute\_pipeline.py }\AttributeTok{{-}{-}project}\NormalTok{ templates/template\_methods\_paper }\AttributeTok{{-}{-}core{-}only}
\end{Highlighting}
\end{Shaded}
\end{frame}

\begin{frame}[fragile,allowframebreaks]{Generated artifact registry}
\protect\phantomsection\label{generated-artifact-registry}
The analysis script writes the following artifacts under
\breaktt{projects/templates/template\_methods\_paper/output/}:

{\def\LTcaptype{none} % do not increment counter
\begin{longtable}[]{@{}
  >{\raggedright\arraybackslash}p{(\linewidth - 2\tabcolsep) * \real{0.5000}}
  >{\raggedright\arraybackslash}p{(\linewidth - 2\tabcolsep) * \real{0.5000}}@{}}
\toprule\noalign{}
\begin{minipage}[b]{\linewidth}\raggedright
Artifact
\end{minipage} & \begin{minipage}[b]{\linewidth}\raggedright
Produced by
\end{minipage} \\
\midrule\noalign{}
\endhead
\breaktt{data/pbspreparation\_worklist.md},
\breaktt{data/pbspreparation\_plan.csv},
\breaktt{data/pbspreparation\_graph.mmd},
\breaktt{data/pbspreparation\_plan.json} & \breaktt{compile\_method()} +
exporters, for \texttt{PBSPreparation} \\
\breaktt{data/sensorcalibrationsweep\_worklist.md},
\breaktt{data/sensorcalibrationsweep\_plan.csv},
\breaktt{data/sensorcalibrationsweep\_graph.mmd},
\breaktt{data/sensorcalibrationsweep\_plan.json} &
\breaktt{compile\_method()} + exporters, for
\breaktt{SensorCalibrationSweep} \\
\breaktt{data/compiled\_plans.json} & Per-method plan summary, consumed
by \breaktt{src/manuscript\_variables.py} \\
\breaktt{reports/gate\_report.json} & \texttt{run\_all\_gates()} tally
across both methods \\
\breaktt{reports/trust\_chain\_report.json} &
\texttt{append\_record()}/\texttt{verify\_chain()} demonstration
chain \\
\breaktt{figures/step\_counts.png} & Step-count bar chart \\
\breaktt{data/manuscript\_variables.json} & Every generated-variable
value, written by \breaktt{z\_generate\_manuscript\_variables.py} \\
\bottomrule\noalign{}
\end{longtable}
}

The \texttt{output/} tree is disposable and regenerated on every run; it
is not the source of truth.
\end{frame}

\begin{frame}[fragile,allowframebreaks]{Determinism}
\protect\phantomsection\label{determinism}
\begin{itemize}
\tightlist
\item
  \breaktt{compile\_method()} is deterministic by construction: the plan
  hash is computed from a canonical, sort-keys JSON payload over
  \texttt{method\_name}, \texttt{method\_version}, \texttt{target}, and
  each scheduled step's identifying fields --- never from a wall-clock
  timestamp or a UUID.
\item
  \breaktt{topological\_order()} breaks scheduling ties by ascending
  \texttt{step\_id}, so the same \texttt{Method} object always yields
  the same step order across processes and platforms.
\item
  Yes --- \breaktt{src/manuscript\_variables.py::generate\_variables}
  recompiles every example method twice at manuscript-build time and
  compares hashes live, so this guarantee is checked on every build, not
  merely asserted once in a test.
\end{itemize}
\end{frame}

\begin{frame}[fragile,allowframebreaks]{Verification (no
hand-transcribed numbers)}
\protect\phantomsection\label{verification-no-hand-transcribed-numbers}
Every quantitative claim in \breakseq{[@sec:results]} is either a generated
variable sourced from a live analysis output or registered in
\breaktt{data/claim\_ledger.yaml} for evidence-registry validation. The
manuscript intentionally does not hand-transcribe volatile values, so
prose and artifacts cannot disagree. Configuration provenance is itself
injected: \breaktt{2107a29383e5e491} is the SHA-256 of
\breaktt{manuscript/config.yaml} at build time, and
\breaktt{2026-06-30T21:41:36Z} records when the variables were generated
(honoring \breaktt{SOURCE\_DATE\_EPOCH} for byte-reproducible builds).
\end{frame}

\end{document}
